% 08-month-by-month. The first half of what used to be one three-figure frame.
% The split is by WHAT THE FIGURE SHOWS, not merely to make room: these two are
% both monthly views of the same 347 paired days — the grid resolves them in
% space, the series in time — and the water-year summary that used to sit with
% them is now 09.
%
% The split also fixes the geometry the old frame fought. Two landscape figures
% stacked need about 0.91x their width in height, which is more than the body
% has, so both had to be set well under the column width and the frame carried
% ~2cm of unusable margin. Pairing the PORTRAIT grid with ONE landscape figure
% lets each take its natural size.
%
% Numbers are from results/FINDINGS.md, "Satellite validation, WY2023",
% converted from fraction to percentage points (-0.043 -> -4.3 pp,
% 0.057 -> 5.7 pp) so the block reads in the same unit as the axes beside it.
% The annual mean is far below the winter monthly values on the plot because
% the 347 days include a summer with no snow on either side.
\begin{frame}[t]
\frametitle{fSCA Bias Month by Month, Water Year 2023}
\vspace{0.15em}
\begin{columns}[c,onlytextwidth]
\begin{column}{0.415\textwidth}
{\centering\includegraphics[height=6.00cm]{water_year_2023_monthly_spatial_grid.png}\par}
\end{column}
\begin{column}{0.570\textwidth}
{\centering\includegraphics[width=0.97\linewidth]{water_year_2023_monthly_bias_mae.png}\par}
\vspace{0.55em}
\begin{block}{MERRA-2 minus MODSCAG}
\scriptsize
\renewcommand{\arraystretch}{1.30}
\begin{tabular}{@{}l@{\hspace{1.1em}}r@{}}
paired days & 347 of 365\\
mean bias   & $-4.3$ pp\\
MAE         & 5.7 pp\\
\end{tabular}
\end{block}
\end{column}
\end{columns}
\vspace{0.30em}
\figcaption{The model runs low against the satellite through accumulation and
converges at melt-out.}
\end{frame}
